\section{The Zermelo-Fraenkel Axioms and the Axiom of Choice}

We begin by reviewing the Zermelo-Fraenkel Axioms of Set Theory. We work in a first-order language with only one relation symbol, denoted $\in$. In principle, we would want to distinguish between the symbol in the language and its interpretation in any model, but we will not do this in practice. In situations where we really want to be pedantic, we will denote the symbol $\dot\in$ and the interpretation $\in$.

We begin with the Axiom of Extensionality.

\subsection{Extensionality}

\begin{baxiom}[The Axiom of Extensionality]\label{ZFC:Ext}
    \begin{align*}
        \forall A \, \forall B \, \brac{\parenth{\forall x \parenth{x \in A \lr x \in B}} \to \parenth{A = B}}
    \end{align*}
\end{baxiom}

We define the $\subseteq$ symbol, used infix as $A \subseteq B$, to be shorthand for the formula $\forall x \, \parenth{x \in A \to x \in B}$. This is not a symbol in the language of set theory; it is merely a convenient abbreviation that we will use hereafter without further discussion.

\Cref{ZFC:Ext} essentially tells us that
\begin{align*}
    \forall A \forall B \parenth{\parenth{A \subseteq B \land B \subseteq A} \to A = B}
\end{align*}
(This is provably equivalent to \Cref{ZFC:Ext}.)

We now define the Axiom Scheme of Comprehension.

\subsection{Comprehension}

First, some notation.

\begin{boxnotation}
    We will denote finite tuples using overlines, so that $\bar{y}$ refers to $\parenth{y_0, \ldots, y_{n-1}}$. We will typically use bars to bundle together variables.
\end{boxnotation}

ZFC is an infinite theory. We call the following an Axiom ``Scheme'' to underscore the fact that it is really an infinite family of axioms bundled together.

\begin{baxiom}[The Axiom (Scheme) of Comprehension]\label{ZFC:Compr}
    Let $x, \bar{y}$ be free variables. For all formulae $\varphi\of{x, \bar{y}}$, the following is an axiom:
    \begin{align*}
        \forall A \forall \bar{y} \exists B \forall x \brac{x \in B \lr \parenth{x \in A \land \varphi\of{x, \bar{y}}}}
    \end{align*}
    Intuitively, this means we can define
    \begin{align*}
        B = \setst{x \in A}{\varphi\of{x, \bar{y}}}
    \end{align*}
    This is an axiom scheme because we can increase the arity of $\varphi$: we don't know how many variables $y_i$ are bundled together in $\bar{y}$, only that for any given statement of the axiom, there are finitely many $y_i$.
\end{baxiom}

Note that \Cref{ZFC:Ext} tells us that we can replace the existential quantifier for $B$ in \Cref{ZFC:Compr} with an existence and uniqueness quantifier without changing its meaning.

Before proceeding further, we define the empty set. In first order logic, structures are required to have non-empty universes. Thus, for every structure $\parenth{A, E}$ in the language of set theory, with $E \subseteq \parenth{A \times A}$ representing equality, we must have
\begin{align*}
    \parenth{A, E} \models \exists x \parenth{x = x}
\end{align*}
The Completeness Theorem then tells us that
\begin{align*}
    \vdash \exists x \parenth{x = x}
\end{align*}
ie, that the formula $\exists x \parenth{x = x}$ \textit{must be a theorem} in the language of sets. Hence, \Cref{ZFC:Ext} and \Cref{ZFC:Compr} tell us that
\begin{align*}
    \vdash \exists! A \, \forall x \, \parenth{x \notin A}
\end{align*}
In other words, it is a theorem in the language of sets that there is a unique set that contains no members whatsoever. This unique set is denoted $\emptyset$, and is called the \textbf{empty set}. Its existence (and uniqueness) is not a distinct axiom, but a direct consequences of \Cref{ZFC:Ext} and \Cref{ZFC:Compr}.

Next, we give the axiom of pairing.

\subsection{Pairing}

\begin{baxiom}[The Axiom of Pairing]\label{ZFC:Pairing}
    \begin{align*}
        \forall x \, \forall y \, \exists A \, \parenth{x \in A \land y \in A}
    \end{align*}
\end{baxiom}

\Cref{ZFC:Pairing} essentially gives us a way of constructing, for any $x$ and $y$, the unique set $A = \set{x, y}$ with the property that
\begin{align*}
    \forall z \parenth{z \in A \lr \parenth{z = x} \lor \parenth{z = y}}
\end{align*}
In particular, it allows us to construct, for any $x$, the set $\set{x, x}$, which, by \Cref{ZFC:Ext}, is exactly $\set{x}$. In other words, it tells us that we can stick any $x$ into a set that only contains $x$.

This allows us to construct pairwise unions, but not arbitrary unions.

\subsection{Unions}

Now, we state the axiom of unions.

\begin{baxiom}[The Axiom of Unions]
    If we denote by $\calF$ a family of sets,
    \begin{align*}
        \forall \calF \, \exists B \, \forall A \, \brac{A \in \calF \to A \subseteq B}
    \end{align*}
\end{baxiom}

As usual, we can use \Cref{ZFC:Ext}, \Cref{ZFC:Compr} to show that
\begin{align*}
    \bigcup \calF = \setst{x}{\exists A \subseteq \calF \, \parenth{x \in A}}
\end{align*}
is well (and uniquely) defined.

Next comes the power set axiom.

\subsection{Power Sets}

\begin{baxiom}[The Axiom of Power Sets]
    \begin{align*}
        \forall A \, \exists \calF \, \forall X \, \brac{X \subseteq A \to X \in \calF}
    \end{align*}
\end{baxiom}

Again, \Cref{ZFC:Ext,ZFC:Compr} tell us that this is equivalent to
\begin{align*}
    \forall A \, \exists! \calF \, \forall X \, \parenth{X \in \calF \lr X \subseteq A}
\end{align*}
This justifies the definition
\begin{align*}
    \Powset(A) := \setst{X}{X \subseteq A}
\end{align*}

Next comes the second (and last) axiom scheme in ZFC, the axiom scheme of replacement.

\subsection{Replacement}

\begin{baxiom}[The Axiom (Scheme) of Replacement]\label{ZFC:Replacement}
    For each formula $\varphi\of{x, y, \bar{z}}$ that does not contain $B$ as a free variable, the following is an axiom:
    \begin{align*}
        \forall A \, \forall \bar{z} \, \brac{
            \brac{\forall x \in A \, \exists! y \, \varphi\of{x, y, \bar{z}}}
            \to
            \brac{\brac{\exists B \, \forall x \in A \, \exists y \in B \, \varphi\of{x, y, \bar{z}}}}
        }
    \end{align*}
\end{baxiom}

% [CLAUDE] draw a diagram for this.

We can unpack this slightly complicated formula using \Cref{ZFC:Ext} and \Cref{ZFC:Compr}: it essentially tells us that we can define
\begin{align*}
    \setst{y}{\exists x \in A \, \parenth{\varphi\of{x, y, \bar{z}}}}
\end{align*}
Ie, if for every $x$, there is a unique witness $y$ for which the formula $\varphi\of{x, y, \bar{z}}$ holds, we can collect all these witnesses into a set that is contained $B$. In particular, ``the number of such witnesses is \textit{small enough} that it \textit{can be collected into a set}.''

Consider the case where $\varphi\of{x, y, \bar{z}}$ defines a function with domain $A$. If we can define $A \times B = \setst{\parenth{x, y}}{x \in A \text{ and } y \in B}$ (which we have yet to show we can do, but let's say we can), then we can put
\begin{align*}
    f = \setst{\parenth{x, y} \in A \times B}{\varphi\of{x, y, z}}
\end{align*}
which is a set by \Cref{ZFC:Compr}. Effectively, \Cref{ZFC:Replacement} tells us that \textit{if the domain of a function is a set, then so is its range}. \Cref{ZFC:Replacement} is an essential step in the process of actually \textit{defining} functions (or showing that this is something we can \textit{do}).

One situation in which \Cref{ZFC:Replacement} is used (and will be used once sufficiently many axioms are stated) is in the situation we talked about at the start of \Cref{Ch1:CH}. We'll see that there is a formula $\varphi\of{n, S}$ so that $\varphi\of{n, S} \lr \parenth{n \in \omega \land S = V_n}$. Then, \Cref{ZFC:Replacement} would imply that
\begin{align*}
    \grpres{V_n}{n < \omega}
\end{align*}
really is a sequence. This allows us to define $V_\omega = \bigcup_{n < \omega} V_n$ and continue.

\subsection{Foundation}

\begin{baxiom}[The Axiom of Foundation]
    \begin{align*}
        \forall A \parenth{A \ne \emptyset \to \exists x \in A \, \forall y \in A \, y \not\in x}
    \end{align*}
\end{baxiom}

We now state the most interesting of the ZFC axioms: the Axiom of Choice.

\subsection{Choice}

\begin{baxiom}[The Axiom of Choice - (AC)]\label{ZFC:AC}
    \hfill
    \begin{align*}
        \forall \calF \, \exists c \, \brac{\forall A \in \calF \, \parenth{A \neq \emptyset \to A \in \dom{c} \land c(A) \in A}}
    \end{align*}
    where $\calF$ is a family of sets, $c$ is a function, and $\dom{c}$ is the domain of $c$. We call $c$ a \textbf{Choice Function on $\calF$}.
\end{baxiom}

The reason we state the Axiom of Choice at this point in our discussion is that modulo the ZFC Axioms that we have seen so far, the Axiom of Choice is equivalent to each of the following.
\begin{enumerate}
    \item $\forall A \, \exists R \subseteq A \times A \, \parenth{A, R}$ is a well-ordering.
    \item $\forall A \, \exists \alpha \, \exists f : A \xrightarrow{\sim} \alpha$, where $\alpha$ is an ordinal.
\end{enumerate}

Note that none of the above axioms imply the existence of the natural numbers, without which virtually none of the mathematics we know and love would be possible. We therefore come to the axiom of infinity, which posits the existence of infinite sets, the `smallest' of which is the natural numbers.

\subsection{Infinity}

\begin{baxiom}[The Axiom of Infinity]\label{ZFC:Infty}
    \hfill
    \begin{align*}
        \exists I \, \parenth{\emptyset \in I \land \forall x \in I \, \parenth{x \cup \set{x} \in I}}
    \end{align*}
    A set $I$ as above is called an \textbf{inductive set}, and it is an axiom of ZFC that such a set exists.
\end{baxiom}

It is possible to prove that given the existence of inductive sets (as posited above), there is a unique inductive set $\omega$ such that for all inductive sets $J$, $\omega \subseteq J$. Indeed, for any inductive set $K$, $\omega$ can be defined as the intersection of all inductive subsets $J \subseteq K$. It can then be shown that this intersection $\omega$ lives in all inductive sets. This minimal inductive set $\omega$ is known as the \textbf{set of natural numbers}. It is easy to see that the famed axioms of Peano are satisfied by $\omega$.

Intuitively, we'd want to say something like
\begin{align*}
    \omega = \bigcap \setst{k}{k \text{ is inductive}}
\end{align*}
The only problem is, a priori, the above intersection is not a set! This is why we need this axiom.
